\documentclass[12pt]{article}
\usepackage[utf8]{inputenc}
\usepackage[catalan]{babel}
\usepackage[margin=1.25in]{geometry} % widen page margins for more whitespace

\begin{document}

\section*{Privadesa demostrable en LLMs amb ZKP: cap a una governança transparent de la IA}



\textbf{Nom de l'estudiant:} Jan Aguiló Plana \\
\textbf{Tutora:} Vanesa Daza

\section*{Objectius del treball}
L’objectiu principal del treball és desenvolupar un sistema que permeti verificar, mitjançant proves de coneixement zero (ZKP), que les respostes generades per models de llenguatge de gran escala (LLMs) no contenen informació confidencial o dades personals. Es vol demostrar que és factible aplicar aquest mecanisme de verificació sense revelar ni el contingut de la resposta ni el funcionament intern del filtre de seguretat. Finalment, es quantificarà el cost real del sistema amb mètriques concretes per a cada consulta (temps de generació i verificació de la prova, ús de memòria, mida de la prova) així com el cost associat al procés inicial de compilació, configuració i generació de claus criptogràfiques.


\section*{Context i motivació}

Un dels grans reptes de la intel·ligència artificial avui, i que només creixerà amb el temps, és la manca de garanties de privadesa. Avui en dia, no és possible assegurar que un model de llenguatge no reveli informació sensible o confidencial en les seves respostes, i això comporta filtracions greus de dades personals, secrets comercials o informació protegida per llei. Aquesta manca de control m’ha obsessionat des que vaig començar a estudiar IA, i aquest treball de final de grau representa per a mi un pas clau per contribuir a una IA més segura, responsable i alineada amb els drets fonamentals de les persones.

Les proves de coneixement zero (zero-knowledge proofs, ZKP) són un tipus de prova criptogràfica que permeten demostrar que una afirmació és certa sense revelar cap informació addicional. En el context d’aquest treball, ens permeten certificar que una resposta generada per un model ha passat un filtre de privadesa (per exemple, que no conté informació personal identificable) sense revelar ni el contingut de la resposta ni els criteris exactes del filtre. Actualment, les tècniques per filtrar sortides d’un model requereixen confiança cega en el sistema o revisió humana, amb els riscos i costos que això comporta. Les ZKP obren la porta a sistemes verificables i automatitzats, amb garanties criptogràfiques, per assegurar que la IA compleix amb criteris de privadesa sense sacrificar la confidencialitat ni la transparència.

\end{document}
